\documentclass{beamer}

\usetheme{Berlin}
\usecolortheme{Orchid}
\useoutertheme{miniframes}
\setbeamertemplate{navigation symbols}{}

\usepackage[utf8]{inputenc}
\usepackage[T1]{fontenc}
\usepackage{amsmath}
\usepackage{amssymb}
\usepackage{booktabs}
\usepackage{graphicx}
\graphicspath{{../images/}}

\usepackage{xcolor}
\usepackage{listings}
\usepackage{hyperref}
\usepackage{tikz}
\usetikzlibrary{arrows.meta,positioning,shapes.geometric,calc}

\lstset{
  language=Java,
  basicstyle=\ttfamily\small,
  keywordstyle=\color{blue},
  commentstyle=\color{gray},
  stringstyle=\color{teal},
  showstringspaces=false,
  columns=fullflexible,
  breaklines=true
}

% Per-slide source line (references shown on the slide itself).
\newcommand{\slidesources}[1]{%
  \vfill
  {\tiny\color{gray}\raggedright\sloppy\parbox{\linewidth}{Sources: #1}\par}%
}

% Unit-wise source strings for this subject (used by slides/notes).
% Path is relative to the lecture `latex/` folder (the compilation CWD).
% OOPS with Java (CSEG1044) — unit-wise sources
%
% These are short, student-facing reference strings intended to be used with:
%   - \slidesources{...}  (in beamer slides)
%   - \notesources{...}   (in lecture notes)
%
% Style inspiration: other subjects in My-Subjects/ use semicolon-separated sources
% and include an "accessed" date for web documentation.

\newcommand{\OOPSUnitISources}{Schildt, \textit{Java: A Beginner's Guide} (10e), McGraw-Hill (Java fundamentals + OOP basics).; Oracle Java Tutorials: Getting Started + Language Basics + Classes and Objects (accessed \today).; Java SE API docs (e.g., \texttt{String}, \texttt{Scanner}) (accessed \today).}

\newcommand{\OOPSUnitIISources}{Schildt, \textit{Java: The Complete Reference} (12e), McGraw-Hill (classes, inheritance, packages, interfaces).; Bloch, \textit{Effective Java} (3e), Addison-Wesley, 2018 (design choices; interfaces vs abstract classes).; Oracle Java Tutorials: Classes and Objects + Interfaces + Packages (accessed \today).; Java SE API docs (e.g., \texttt{Comparable}, \texttt{Runnable}) (accessed \today).}

\newcommand{\OOPSUnitIIISources}{Schildt, \textit{Java: The Complete Reference} (12e) (nested/inner classes, exceptions, threads, I/O).; Oracle Java Tutorials: Nested Classes + Exceptions + Concurrency + Basic I/O (accessed \today).; Java SE API docs (e.g., \texttt{Thread}, \texttt{AutoCloseable}) (accessed \today).}

\newcommand{\OOPSUnitIVSources}{Schildt, \textit{Java: The Complete Reference} (12e) (generics, lambdas, Swing, JDBC).; Bloch, \textit{Effective Java} (3e) (generics + lambdas).; Oracle Java Tutorials: Generics + Lambda Expressions + Swing + JDBC Basics (accessed \today).; Java SE API docs (e.g., \texttt{java.util.function}, \texttt{javax.swing}, \texttt{java.sql}) (accessed \today).}

\newcommand{\OOPSUnitVSources}{Schildt, \textit{Java: The Complete Reference} (12e) (Collections Framework + wrapper classes).; Oracle Java docs: Collections Framework overview (accessed \today).; Java SE API docs (\texttt{java.util} collections; wrapper classes in \texttt{java.lang}) (accessed \today).}

\newcommand{\OOPSUnitVISources}{Git SCM: \textit{Pro Git} (2e) (version control workflow), accessed \today.; GitHub Docs (repositories, issues, pull requests), accessed \today.; Oracle Java Tutorials: Swing + JDBC (accessed \today).; JUnit 5 User Guide (accessed \today).; Apache Maven Project: Getting Started (accessed \today).; Gradle User Manual: Getting Started (accessed \today).; UML class diagrams: UML 2.x overview/spec (accessed \today).}

\newcommand{\OOPSMiscWhyInterfacesSources}{Schildt, \textit{Java: The Complete Reference} (12e) (interfaces).; Bloch, \textit{Effective Java} (3e) (prefer interfaces where appropriate).; Oracle Java Tutorials: Interfaces (accessed \today).; Java SE API docs (e.g., \texttt{Comparable}, \texttt{Runnable}) (accessed \today).}



\title[OOPS with Java]{Object Oriented Programming (CSEG1044)}
\subtitle{Unit 04 -- Lecture 01: Generics Fundamentals (Generic Classes)}
\begin{document}

\begin{frame}[fragile]
  \titlepage
  \vspace{0.5em}
  \slidesources{\OOPSUnitIVSources}
\end{frame}

\begin{frame}{Quick Links}
  \centering
  \hyperlink{sec:overview}{\beamerbutton{Overview}}\hspace{0.6em}
  \hyperlink{sec:concepts}{\beamerbutton{Key Topics}}\hspace{0.6em}
  \hyperlink{sec:demo}{\beamerbutton{Demo}}\hspace{0.6em}
  \hyperlink{sec:summary}{\beamerbutton{Summary}}
  \slidesources{\OOPSUnitIVSources}
\end{frame}

\begin{frame}{Agenda}
  \tableofcontents
  \slidesources{\OOPSUnitIVSources}
\end{frame}

\section{Overview}
\label{sec:overview}

\begin{frame}{Lecture Snapshot}
\begin{itemize}[<+->]
  \item \textbf{Unit focus:} Generics, Lambdas, Swing, and JDBC
  \item \textbf{Course thread:} Use Generics Fundamentals (Generic Classes) to move Campus Resource Manager toward a real desktop application with UI, callbacks, and persistence.
\end{itemize}
  \slidesources{\OOPSUnitIVSources}
\end{frame}

\begin{frame}{Recall Before You Start}
\begin{itemize}[<+->]
  \item \textbf{Classes and methods}: A generic class is still a class; generics only add type parameters to familiar structures. Main reference: \texttt{Unit 02 / Lecture 01 slides/notes}.
  \item \textbf{Type relationships}: Generic type safety is easier to read once basic inheritance intuition is already clear. Main reference: \texttt{Unit 02 / Lecture 04 slides/notes}.
\end{itemize}
  \slidesources{\OOPSUnitIVSources}
\end{frame}


\begin{frame}{Learning Outcomes}
  \begin{itemize}[<+->]
    \item Explain why generic classes are better than duplicated classes or \texttt{Object}-only designs.
    \item Read and write a generic class such as \texttt{ResourceHolder<T>}.
    \item Distinguish \textbf{type parameter}, \textbf{type argument}, and \textbf{parameterized type}.
    \item Identify raw-type usage and explain why it weakens type safety.
  \end{itemize}
  \slidesources{\OOPSUnitIVSources}
\end{frame}

\begin{frame}{Study Flow for This Lecture}
\begin{enumerate}[<+->]
  \item Refresh the prerequisite bridge before reading new ideas in isolation.
  \item Study the main build in this order: \textit{Why this topic exists} $\rightarrow$ \textit{Attempt 1: separate classes for every type} $\rightarrow$ \textit{Attempt 2: use one class with \texttt{Object}} $\rightarrow$ \textit{Why \texttt{Object} is unsafe}.
  \item Trace the worked example and connect each line back to one concept frame.
  \item Attempt the mini exercises before moving to the recap and exit check.
\end{enumerate}
  \slidesources{\OOPSUnitIVSources}
\end{frame}


\section{Key Topics}
\label{sec:concepts}

\begin{frame}{Unit 04: Generics, Lambdas, GUI Swing \& Database Connectivity}
  \begin{itemize}[<+->]
    \item Lecture focus: Generics Fundamentals (Generic Classes)
  \end{itemize}
  \slidesources{\OOPSUnitIVSources}
\end{frame}

\begin{frame}{Today: Teaching Path}
  \begin{enumerate}[<+->]
    \item Same problem solved with many separate classes
    \item Same problem solved with one \texttt{Object}-based class
    \item Same problem solved with one generic class
    \item Correct usage, raw-type mistakes, and a small project connection
  \end{enumerate}
  \slidesources{\OOPSUnitIVSources}
\end{frame}

\begin{frame}{Start with the problem}
  \begin{itemize}[<+->]
    \item Suppose the project stores different campus resources: books, laptops, and lab keys.
    \item A beginner often creates a separate holder class for each type.
    \item That works, but the class logic becomes repeated with only the data type changing.
  \end{itemize}
  \slidesources{\OOPSUnitIVSources}
\end{frame}

\begin{frame}[fragile]{Repeated classes: works, but duplicates logic}
\begin{columns}[T]
\column{0.48\textwidth}
\begin{lstlisting}
class BookHolder {
    private Book item;
    void set(Book item) {
        this.item = item;
    }
    Book get() {
        return item;
    }
}
\end{lstlisting}
\column{0.48\textwidth}
\begin{lstlisting}
class LaptopHolder {
    private Laptop item;
    void set(Laptop item) {
        this.item = item;
    }
    Laptop get() {
        return item;
    }
}
\end{lstlisting}
\end{columns}
\vspace{0.4em}
\begin{itemize}[<+->]
  \item Structure is the same; only the type changes.
\end{itemize}
  \slidesources{\OOPSUnitIVSources}
\end{frame}

\begin{frame}{The first improvement many students try}
  \begin{itemize}[<+->]
    \item ``Why not store everything as \texttt{Object}?'' 
    \item This removes duplication because one class can hold any reference type.
    \item But now the compiler no longer knows the real type stored inside.
  \end{itemize}
  \slidesources{\OOPSUnitIVSources}
\end{frame}

\begin{frame}[fragile]{One \texttt{Object}-based class: less duplication, more risk}
\begin{lstlisting}
class ResourceHolder {
    private Object item;

    void set(Object item) { this.item = item; }
    Object get() { return item; }
}
\end{lstlisting}
\begin{itemize}[<+->]
  \item You can store a \texttt{Book}, \texttt{Laptop}, or \texttt{LabKey}.
  \item But \texttt{get()} returns \texttt{Object}, so the caller must cast.
\end{itemize}
  \slidesources{\OOPSUnitIVSources}
\end{frame}

\begin{frame}[fragile]{Why \texttt{Object} becomes unsafe}
\begin{lstlisting}
ResourceHolder holder = new ResourceHolder();
holder.set(new Laptop("L-14"));

Book b = (Book) holder.get();   // compiles
\end{lstlisting}
\begin{itemize}[<+->]
  \item The cast compiles.
  \item The data is actually a \texttt{Laptop}, not a \texttt{Book}.
  \item Error appears at runtime as \texttt{ClassCastException}.
\end{itemize}
  \slidesources{\OOPSUnitIVSources}
\end{frame}

\begin{frame}{Generics solve both sides of the problem}
\begin{itemize}[<+->]
  \item We still want \textbf{one reusable class}.
  \item We also want \textbf{compile-time type checking}.
  \item A generic class gives both: reuse \textit{and} safety.
\end{itemize}
  \slidesources{\OOPSUnitIVSources}
\end{frame}

\begin{frame}[fragile]{Generic class design}
\begin{lstlisting}
class ResourceHolder<T> {
    private T item;

    void set(T item) { this.item = item; }
    T get() { return item; }
}
\end{lstlisting}
\begin{itemize}[<+->]
  \item \texttt{T} is a type parameter.
  \item Inside the class, \texttt{T} behaves like ``the type we will decide later''.
\end{itemize}
  \slidesources{\OOPSUnitIVSources}
\end{frame}

\begin{frame}{How to read the syntax}
\begin{itemize}[<+->]
  \item \texttt{class ResourceHolder<T>} means ``define a class that works for some type \texttt{T}''.
  \item \texttt{ResourceHolder<Book>} means ``this particular holder stores \texttt{Book} objects''.
  \item \texttt{T} is the \textbf{type parameter}; \texttt{Book} is the \textbf{type argument}.
  \item After substitution, \texttt{ResourceHolder<Book>} is a \textbf{parameterized type}.
\end{itemize}
  \slidesources{\OOPSUnitIVSources}
\end{frame}

\begin{frame}[fragile]{Using the same class with different resource types}
\begin{lstlisting}
ResourceHolder<Book> bookHolder = new ResourceHolder<>();
ResourceHolder<Laptop> laptopHolder = new ResourceHolder<>();
ResourceHolder<LabKey> keyHolder = new ResourceHolder<>();
\end{lstlisting}
\begin{itemize}[<+->]
  \item Same class definition
  \item Different type-safe usages
  \item No repeated holder classes
\end{itemize}
  \slidesources{\OOPSUnitIVSources}
\end{frame}

\begin{frame}{What the compiler now protects}
\begin{itemize}[<+->]
  \item A \texttt{ResourceHolder<Book>} accepts \texttt{Book} values, not \texttt{Laptop} values.
  \item \texttt{get()} returns \texttt{Book}, so no cast is needed.
  \item Many mistakes move from runtime to compile time.
\end{itemize}
  \slidesources{\OOPSUnitIVSources}
\end{frame}

\begin{frame}{Important rule: primitives are not type arguments}
\begin{itemize}[<+->]
  \item \texttt{ResourceHolder<int>} is invalid.
  \item Generics need reference types.
  \item Use wrapper classes such as \texttt{Integer}, \texttt{Double}, and \texttt{Character}.
\end{itemize}
  \slidesources{\OOPSUnitIVSources}
\end{frame}

\begin{frame}{Raw types: old style, unsafe style}
\begin{itemize}[<+->]
  \item Raw type means writing \texttt{ResourceHolder} instead of \texttt{ResourceHolder<Book>} or another parameterized type.
  \item Java allows it for backward compatibility with pre-generics code.
  \item But you lose most compile-time type checking.
  \item In modern code, warnings about raw types should usually be treated as design problems.
\end{itemize}
  \slidesources{\OOPSUnitIVSources}
\end{frame}

\begin{frame}{Where you will meet generic classes again}
\begin{itemize}[<+->]
  \item Java library classes use the same idea: collections, maps, queues, optional wrappers, and more.
  \item In your own project, generic classes are useful for reusable containers, response wrappers, repositories, and helper utilities.
  \item This lecture only builds the foundation; the next lecture extends it to methods, bounds, and wildcards.
\end{itemize}
  \slidesources{\OOPSUnitIVSources}
\end{frame}

\begin{frame}{Checkpoint and Reflection}
\begin{block}{Reflection prompt}
Where would \textit{Generics Fundamentals (Generic Classes)} matter most in Campus Resource Manager, and which design choice would you justify using this lecture's ideas?
\end{block}
\begin{block}{Quick self-check}
Which part needs one more example before you move on: \textit{Why this topic exists}, \textit{Attempt 1: separate classes for every type}, the worked example, or the recap?
\end{block}
  \slidesources{\OOPSUnitIVSources}
\end{frame}

\section{Demo / Example}
\label{sec:demo}

\begin{frame}[fragile]{Worked example: sample domain classes}
\begin{lstlisting}[basicstyle=\ttfamily\small]
class Book {
    String title;
    Book(String title) { this.title = title; }
}

class LabKey {
    String roomCode;
    LabKey(String roomCode) { this.roomCode = roomCode; }
}
\end{lstlisting}
  \slidesources{\OOPSUnitIVSources}
\end{frame}

\begin{frame}[fragile]{Worked example: generic holder}
\begin{lstlisting}[basicstyle=\ttfamily\small]
class ResourceHolder<T> {
    private final T item;
    ResourceHolder(T item) { this.item = item; }
    T get() { return item; }
}
\end{lstlisting}
  \slidesources{\OOPSUnitIVSources}
\end{frame}

\begin{frame}[fragile]{Safe usage: one class, different types}
\begin{lstlisting}[basicstyle=\ttfamily\small]
public class DemoSafe {
    public static void main(String[] args) {
        ResourceHolder<Book> b =
            new ResourceHolder<>(new Book("Clean Code"));
        ResourceHolder<LabKey> k =
            new ResourceHolder<>(new LabKey("LAB-2"));

        Book book = b.get();
        LabKey key = k.get();
        System.out.println(book.title + " / " + key.roomCode);
    }
}
\end{lstlisting}
  \slidesources{\OOPSUnitIVSources}
\end{frame}

\begin{frame}[fragile]{Unsafe usage: raw type removes the benefit}
\begin{lstlisting}
public class DemoRaw {
    public static void main(String[] args) {
        ResourceHolder raw =
            new ResourceHolder(new LabKey("LAB-2"));

        Book wrong = (Book) raw.get(); // runtime risk
        System.out.println(wrong.title);
    }
}
\end{lstlisting}
  \slidesources{\OOPSUnitIVSources}
\end{frame}

\begin{frame}{Mini Exercises (2--3 mins)}
  \begin{enumerate}[<+->]
    \item Why is a single \texttt{Object}-based class better than duplicated classes, but still not the best final design?
    \item What is the difference between \texttt{T} in \texttt{ResourceHolder<T>} and \texttt{Book} in \texttt{ResourceHolder<Book>}?
    \item Why do we write \texttt{ResourceHolder<Integer>} instead of \texttt{ResourceHolder<int>}?
  \end{enumerate}
  \slidesources{\OOPSUnitIVSources}
\end{frame}

\begin{frame}{Watch-outs}
\begin{itemize}[<+->]
  \item Treating generic type parameters as if they were inheritance declarations.
  \item Falling back to raw types and losing type safety.
  \item Expecting primitive types to work directly as generic arguments.
\end{itemize}
  \slidesources{\OOPSUnitIVSources}
\end{frame}

\section{Summary}
\label{sec:summary}

\begin{frame}{Key Takeaways}
  \begin{itemize}[<+->]
    \item Separate classes for each type cause duplication.
    \item A single \texttt{Object}-based class removes duplication but pushes errors to runtime.
    \item A generic class such as \texttt{ResourceHolder<T>} gives reuse with compile-time checking.
    \item Raw types weaken the protection that generics are meant to provide.
  \end{itemize}
  \slidesources{\OOPSUnitIVSources}
\end{frame}

\begin{frame}{Checkpoint Questions}
  \begin{enumerate}[<+->]
    \item What design problem leads us toward generic classes?
    \item What is the meaning of \texttt{ResourceHolder<Book>}?
    \item Why is raw \texttt{ResourceHolder} unsafe compared with \texttt{ResourceHolder<Book>}?
  \end{enumerate}
  \slidesources{\OOPSUnitIVSources}
\end{frame}

\begin{frame}{Next Study Steps}
\begin{itemize}[<+->]
  \item Revisit the recall references if any older idea still feels weak.
  \item Rework the lecture example without looking at the notes first.
  \item Attempt the self-study practice and then answer the exit questions in full sentences.
  \item Apply the idea once inside Campus Resource Manager to make the concept stick.
\end{itemize}
  \slidesources{\OOPSUnitIVSources}
\end{frame}

\begin{frame}{Next Lecture Preview}
  \textbf{Next:} Generic methods, bounded type parameters, and wildcards
  \vspace{0.6em}
  \begin{itemize}[<+->]
    \item We will move from generic classes to methods and then learn how to restrict acceptable types.
  \end{itemize}
  \slidesources{\OOPSUnitIVSources}
\end{frame}

\end{document}
